\chapter{Registros da Dinâmica do Grupo de Testes}

Este apêndice apresenta, de forma simples e direta, alguns registros textuais das interações do grupo de testes sobre o uso do celular. Os trechos a seguir foram mantidos próximos da forma original, de modo a ilustrar percepções, dificuldades e estratégias relatadas pelos participantes.

\begin{quote}
Uma das minhas metas vai ser parar de usar o celular na hora do treino da academia. Perco um tempão e não rendo direito porque fico olhando o Instagram toda hora.
\end{quote}

\begin{quote}
O meu vício é com aqueles vídeos curtos de meme, fico ali vidrado e quando vi já se passaram horas.
\end{quote}

\begin{quote}
Às vezes fico mexendo no celular antes de dormir também.
\end{quote}

\begin{quote}
Enviando diariamente, da pra notar como é a sensação de ter outras pessoas vendo seu tempo impacta nisso.
\end{quote}

\begin{quote}
Isso acaba com meu sono.
\end{quote}

\begin{quote}
Agora estou tentando ler antes de dormir, isso dá uma ajudada, redes sociais têm muita informação.
\end{quote}

\begin{quote}
A meta diária de tempo de uso saudável é qual? Pelo que pesquisamos aqui, até 3 horas por dia é um limite saudável para uso do celular.
\end{quote}

\begin{quote}
Recentemente eu desativei o Instagram temporariamente, por muitos meses, quase 1 ano. Eu perdia muito tempo, minha vida era sugada.
\end{quote}

\begin{quote}
Uma parada que fiz e me ajudou foi desativar as notificações de todas as redes sociais. Ainda tenho o hábito de abrir Instagram e WhatsApp de tempos em tempos, mas agora a ansiedade de ver não atrapalha nas atividades de rotina da vida.
\end{quote}

\begin{quote}
Tive 638 notificações no dia de ontem, agora que eu vi.
\end{quote}

\begin{quote}
Eu desativei várias notificações no meu celular, principalmente do Instagram. Acho que vou fazer isso, fiquei muito surpreso com esse número.
\end{quote}

\begin{quote}
Fiz algumas mudanças hoje e mais tarde mostro o resultado. Vamos ver se vai dar resultado.
\end{quote}

\begin{quote}
Como eu não tenho computador, uso muito o celular pra tudo, mas ainda assim às vezes preciso maneirar nas redes sociais.
\end{quote}

\begin{quote}
Meu tempo de uso hoje: fiz algumas mudanças e isso reduziu muito o meu tempo no celular, ontem deu 10 horas. As mudanças que fiz foram: não pegar no celular ao acordar, não usar o telefone na academia (eu coloco o fone, escolho a playlist e não mexo mais), não usar o celular enquanto estou fazendo outras tarefas, como comer, arrumar a casa etc. E, principalmente, desativei todas as notificações das redes sociais, deixei apenas o WhatsApp. São mudanças simples, mas fizeram diferença pra mim. Também coloquei limite de tempo de 3 horas pra usar as redes sociais. Com o tempo eu pretendo diminuir.
\end{quote}

\begin{quote}
Algumas mudanças simples já fazem muita diferença. Eu mesmo já desativei todas as notificações há muito tempo. Só vejo se eu entrar no app.
\end{quote}

\begin{quote}
Feriado, né: não fui pra academia e nem pra faculdade, então isso já explica horas e horas. Passei muito tempo de papo no WhatsApp, e evitando o Instagram porque meus dados móveis estão acabando. Além disso, passei parte da tarde e da noite fazendo entregas.
\end{quote}

\begin{quote}
Te entendo totalmente. Daqui a pouco vou mandar a minha aqui também, fica bem grande o tempo de ontem.
\end{quote}

\begin{quote}
Eu e outra pessoa estamos viajando, então acabamos esquecendo de postar aqui. Vamos começar. É pra colocar tempo diário ou média semanal? Tempo diário.

Estamos enviando todo dia até terça da semana que vem, para ver como lidamos com a diminuição e melhoria do tempo de uso do celular em grupo por 1 semana. Estamos nos esforçando e conseguindo diminuir o tempo de uso.
\end{quote}

\begin{quote}
Ontem eu extrapolei um pouco, fiquei boa parte do dia mexendo no celular. Hoje eu quis ficar mais desconectado, aí saí de casa pra passear, viver um pouco. As 5 horas de ontem no celular me assustaram, eu estava perdendo muito tempo no celular.
\end{quote}

\begin{quote}
Meu tempo de hoje foi bem curtinho. Mas também acordei meio tarde, fiquei umas 2 horas na academia e depois rodei o bairro todo atrás de uma peça da máquina de lavar. Nem tive muito tempo pra mexer no celular hoje.
\end{quote}

\begin{quote}
O meu foi horrível. Hoje foi um dia extremamente improdutivo. Acho que o maior desafio é encontrar lazer além das redes sociais. Mas é isso, não dá pra ser perfeito todos os dias. Amanhã eu tento fazer melhor.
\end{quote}

\begin{quote}
Hoje eu usei muito de dia, agora à noite estou fazendo entrega de aplicativo. Usou pra trabalhar, né.
\end{quote}

\begin{quote}
O meu foi terrível de novo, mas amanhã volto à rotina normal de academia, faculdade e trabalho, e acredito que vai dar diferença.
\end{quote}

\begin{quote}
Hoje foi um dia que eu mexi pouco no celular, mas em compensação acabei ficando bastante no computador.
\end{quote}

\begin{quote}
Mexi pouco também e geralmente fico mais no computador, focado em trabalho, faculdade e estudos.
\end{quote}

\begin{quote}
Percebi que a vida de CLT faz eu mexer menos no celular. Nesse sentido, realmente a vida de CLT é boa demais pra não mexer tanto no celular.
\end{quote}

\begin{quote}
Ontem foi terrível porque eu fiquei praticamente o dia em casa sem fazer nada relevante. Hoje eu voltei à programação normal e percebi que, quando eu ocupo a cabeça com outras coisas, mexo menos no celular.
\end{quote}

Esses registros ilustram como, ao longo da dinâmica em grupo, os participantes foram tomando consciência do próprio tempo de tela, experimentando ajustes na rotina e percebendo relações entre uso do celular, sono, produtividade e bem-estar.